\documentclass[aspectratio=169,11pt]{beamer}

\usetheme{Madrid}
\usecolortheme{default}
\usefonttheme{professionalfonts}
\setbeamertemplate{navigation symbols}{}
\setbeamertemplate{footline}[frame number]

\usepackage{amsmath, amssymb, booktabs, graphicx}

\title{Informality, Dualism, and Contract Enforcement}
\subtitle{Labor Markets and Political/Cultural Institutions --- Week 3}
\author{Teaching deck draft}
\date{}

\begin{document}

\begin{frame}
  \titlepage
\end{frame}

\begin{frame}{Central question}
\begin{itemize}
    \item How are labor markets organized when formal contracts, payroll systems, and legal protections cover only part of work?
    \item Which margin is informal: the worker, the firm, the job, payroll reporting, or effective enforceability?
    \item How do worker and firm decisions jointly shape wages, benefits, mobility, productivity, and welfare?
\end{itemize}
\end{frame}

\begin{frame}{Bridge from Week 2}
\begin{itemize}
    \item Week 2: law matters when enforcement, knowledge, and compliance make it effective.
    \item Week 3: partial enforcement creates formal, informal, and partly formal labor exchange.
    \item The formal--informal margin is not one binary variable.
    \item Keep separating law, enforcement, compliance, incidence, and contract quality.
\end{itemize}
\end{frame}

\begin{frame}{What counts as informality?}
\small
\begin{tabular}{p{0.31\linewidth} p{0.55\linewidth}}
\toprule
Object & Labor-market interpretation \\
\midrule
Unregistered firm & business is outside registries, tax systems, or licenses \\
Off-the-books worker & registered firm hides some employment relationships \\
Informal wage job & wage work lacks payroll registration, benefits, or protection \\
Own-account work & self-employment outside formal tax and benefit systems \\
Missing compliance & payroll, wages, benefits, hours, or severance are underreported \\
Unenforceable contract & written or verbal promise cannot be credibly claimed \\
\bottomrule
\end{tabular}
\end{frame}

\begin{frame}{Extensive and intensive margins}
\begin{itemize}
    \item Extensive firm margin: whether the firm registers at all.
    \item Extensive job margin: whether a worker-firm match is reported as formal.
    \item Intensive payroll margin: whether wages, hours, and contributions are fully reported.
    \item Intensive enforcement margin: whether benefits and contract terms are actually claimable.
    \item A registered firm can still employ informal workers.
\end{itemize}
\end{frame}

\begin{frame}{Dualism vs. equilibrium views}
\small
\begin{tabular}{p{0.24\linewidth} p{0.31\linewidth} p{0.31\linewidth}}
\toprule
View & Worker implication & Firm implication \\
\midrule
Dualist & workers rationed out of better formal jobs & informal firms are persistently low productivity \\
Competitive & workers sort by costs, benefits, flexibility, and outside options & firms choose status by taxes, regulation, and enforcement \\
Modern equilibrium & segmentation and sorting can coexist & heterogeneous firms exploit several informal margins \\
Institutional & protection depends on enforceability and state capacity & compliance depends on visibility, sanctions, and formal benefits \\
\bottomrule
\end{tabular}
\end{frame}

\begin{frame}{Anchor papers}
\small
\begin{tabular}{p{0.32\linewidth} p{0.49\linewidth}}
\toprule
Anchor & Use in Week 3 \\
\midrule
Ulyssea overview & map causes, consequences, and policy margins \\
La Porta--Shleifer & dualist productivity-gap benchmark \\
Ulyssea Brazil & firm registration, payroll informality, and structural policy counterfactuals \\
Meghir--Narita--Robin & wages, search, sorting, and compensating differentials \\
Naidu--Yuchtman & contract enforcement as a labor-market institution \\
Samaniego--Fernandez & modern enforcement change in informal employment costs \\
Gerard & social-program incidence around the formal--informal margin \\
\bottomrule
\end{tabular}
\end{frame}

\begin{frame}{Worker-side margins}
\begin{itemize}
    \item Wages: penalties, premia, wage offsets, and compensating differentials.
    \item Benefits: social insurance, health, pensions, leave, severance, injury protection.
    \item Risk: volatility, nonpayment, dismissal, injury, and weak claims.
    \item Mobility: transitions between informality, formality, unemployment, and self-employment.
    \item Sorting: skill, outside options, networks, household constraints, and benefit valuation.
\end{itemize}
\end{frame}

\begin{frame}{Wage gaps are not enough}
\begin{itemize}
    \item A formal--informal wage gap may reflect productivity sorting.
    \item It may reflect rationing out of formal jobs.
    \item It may reflect missing benefits and risk compensation.
    \item It may reflect firm-level evasion and payroll costs.
    \item Interpretation needs transitions, benefits, firm productivity, and worker characteristics.
\end{itemize}
\end{frame}

\begin{frame}{Firm-side margins}
\begin{itemize}
    \item Registration: whether the business becomes visible to the state.
    \item Payroll reporting: whether workers, wages, hours, and contributions are recorded.
    \item Compliance: whether mandated benefits and contract terms are actually paid.
    \item Productivity and scale: formal status can affect credit, buyers, courts, and size.
    \item Adaptation: firms may shrink, outsource, split entities, reclassify workers, or substitute technology.
\end{itemize}
\end{frame}

\begin{frame}{Registration is not payroll compliance}
\small
\begin{tabular}{p{0.24\linewidth} p{0.30\linewidth} p{0.32\linewidth}}
\toprule
State & What is visible? & What can still be hidden? \\
\midrule
Unregistered firm & little or no firm record & workers, wages, output, benefits \\
Registered firm & firm identity, sector, location & off-the-books workers and underreported wages \\
Payroll record & worker match and some earnings & hours, true wage, benefits, enforceability \\
Formal contract & written terms & actual claimability and retaliation risk \\
\bottomrule
\end{tabular}
\end{frame}

\begin{frame}{Contract enforcement}
\[
V^{formal}_{ij}-V^{informal}_{ij}
=B_{ij}-C_{ij}(R,E,\tau)+X_{ij}(K,Q)
\]
\begin{itemize}
    \item $B_{ij}$: benefits of formal status for the worker-firm match.
    \item $C_{ij}$: registration, payroll, tax, compliance, and expected enforcement costs.
    \item $X_{ij}$: gains from contract quality, courts, credit, reputation, or stable matches.
    \item Teaching device: formal status is valuable only through concrete worker and firm margins.
\end{itemize}
\end{frame}

\begin{frame}{Formal contract vs. effective enforceability}
\begin{itemize}
    \item A written contract can exist without accessible claims.
    \item Payroll registration can exist without complete benefit receipt.
    \item Informal relationships can be supported by reputation or repeated interaction.
    \item Naidu--Yuchtman: legal enforceability of labor contracts changed wages and labor-market adjustment.
    \item The key question is whose claims are enforceable and at what cost.
\end{itemize}
\end{frame}

\begin{frame}{Equilibrium interactions}
\begin{itemize}
    \item Workers value cash wages, benefits, risk, flexibility, and mobility differently.
    \item Firms value payroll savings, formal market access, credit, courts, and enforcement avoidance differently.
    \item Stronger enforcement can formalize jobs, shrink firms, change wages, or reallocate employment.
    \item Social programs can change the value of formal or informal status.
    \item The welfare object is a joint worker--firm and policy-incidence problem.
\end{itemize}
\end{frame}

\begin{frame}{Empirical designs and what they identify}
\small
\begin{tabular}{p{0.34\linewidth} p{0.25\linewidth} p{0.27\linewidth}}
\toprule
Design & Identifies best & Observed margin \\
\midrule
Worker panels & sorting and job ladders & transitions, wages, benefits \\
Inspection intensity & enforcement and compliance & payroll, formality, wages, firm response \\
Matched worker--firm models & wage-setting and sorting & wages, flows, firm productivity \\
Structural firm models & registration and payroll choices & scale, evasion, welfare \\
Policy discontinuities & benefit or tax wedges & take-up, formality, welfare \\
Historical legal designs & contract enforceability & wages, mobility, bargaining \\
\bottomrule
\end{tabular}
\end{frame}

\begin{frame}{Diagnostic rule}
\begin{itemize}
    \item Do not say ``informality changes'' without naming the margin.
    \item Did firm registration rise?
    \item Did payroll reporting rise?
    \item Did benefit receipt or contract claimability rise?
    \item Did workers move, wages change, or firms reallocate employment?
    \item The identifying variation determines the object identified.
\end{itemize}
\end{frame}

\begin{frame}{Week 3 lab}
\begin{itemize}
    \item Reproduce: synthetic Ulyssea-style firm margins for registration and payroll reporting.
    \item Diagnose: classify worker-transition and wage patterns using Meghir--Narita--Robin.
    \item Transfer: apply the classification to a modern enforcement setting.
    \item Core output: what side moves, what margin is observed, and what the design identifies.
\end{itemize}
\end{frame}

\begin{frame}{Bridge to Week 4}
\begin{itemize}
    \item Week 2: when does labor law become effective?
    \item Week 3: what happens when formal contracts cover only part of work?
    \item Week 4: how do workers act collectively when individual claims are weak?
    \item Worker voice changes bargaining, monitoring, compliance, and surplus division.
\end{itemize}
\end{frame}

\begin{frame}{Takeaways}
\begin{itemize}
    \item Informality is a two-sided labor-market institution.
    \item Worker formality, firm registration, payroll reporting, and enforceability are different margins.
    \item Dualism and sorting are not mutually exclusive in equilibrium.
    \item Strong empirical work names the observed margin and the identifying variation.
\end{itemize}
\end{frame}

\end{document}
